\section{Sinh Nước Đi}

\begin{frame}{Pipeline sinh nước đi}
  \centering
  \begin{fitdiagram}
    \begin{tikzpicture}[node distance=0.58cm]
      \node[flowbox, text width=3.3cm] (call) {\code{generate\_moves}\\\code{(pos, type)}};
      \node[flowbox, text width=3.3cm, below=of call] (pseudo) {Sinh pseudo-legal\\cho từng quân};
      \node[flowbox, text width=3.3cm, below=of pseudo] (legal) {Lọc \code{pos.legal(m)}};
      \node[flowbox, text width=3.3cm, below=of legal] (type) {Lọc Legal/Capture/Quiet};
      \node[redbox, text width=3.3cm, below=of type] (out) {$\actions(s)$};
      \draw[arrow] (call) -- (pseudo);
      \draw[arrow] (pseudo) -- (legal);
      \draw[arrow] (legal) -- (type);
      \draw[arrow] (type) -- (out);
    \end{tikzpicture}
  \end{fitdiagram}

  \vspace{0.6em}
  \small \code{MoveList} dùng \code{ArrayVec<Move, 128>}: tránh cấp phát heap trong search.
\end{frame}

\begin{frame}{Nhóm quân và kỹ thuật sinh move}
  \small
  \begin{tabular}{p{0.22\textwidth}p{0.28\textwidth}p{0.40\textwidth}}
    \toprule
    Nhóm & Quân & Kỹ thuật \\
    \midrule
    Không bị cản & Tướng, Sĩ, Tốt & Lookup table tĩnh, mask quân nhà. \\
    Nhảy có blocker & Mã, Tượng & Magic Bitboards cho chân/mắt. \\
    Trượt đường thẳng & Xe, Pháo & Rank/File table; Pháo tách quiet và capture. \\
    Kiểm tra chiếu & Tướng, Tốt, Mã, Xe, Pháo & Backward attack scanner. \\
    \bottomrule
  \end{tabular}
\end{frame}

\begin{frame}{Mã bị cản chân}
  \begin{columns}[T,totalwidth=\textwidth]
    \begin{column}{0.42\textwidth}
      \begin{itemize}
        \item Mã có tối đa 8 nước.
        \item Nước nhảy phụ thuộc 4 ô chân Mã.
        \item Blocker mask có $2^4=16$ trạng thái.
      \end{itemize}
      \begin{alertblock}{Ý tưởng}
        Occupancy quanh Mã ánh xạ sang attack table.
      \end{alertblock}
    \end{column}
    \begin{column}{0.54\textwidth}
      \centering
      \includegraphics[height=0.68\textheight]{\figpath{figure-02a-horse-leg.png}}
      \figcaption{3.a}{Mã bị cản chân.}
    \end{column}
  \end{columns}
\end{frame}

\begin{frame}{Magic Bitboards cho Mã}
  \begin{columns}[T,totalwidth=\textwidth]
    \begin{column}{0.45\textwidth}
      \[
        index=\frac{(occupied \& mask)\times magic}{2^{128-bits}}
      \]
      \begin{itemize}
        \item \code{mask}: các ô blocker liên quan.
        \item \code{magic}: hằng số để nén pattern.
        \item \code{table[index]}: tập ô đi được.
      \end{itemize}
    \end{column}
    \begin{column}{0.51\textwidth}
      \centering
      \begin{fitdiagram}
        \begin{tikzpicture}[node distance=0.6cm]
          \node[flowbox] (occ) {occupied};
          \node[flowbox, below=of occ] (mask) {\code{\& mask}};
          \node[flowbox, below=of mask] (mul) {$\times magic$};
          \node[redbox, below=of mul] (idx) {table index};
          \draw[arrow] (occ) -- (mask);
          \draw[arrow] (mask) -- (mul);
          \draw[arrow] (mul) -- (idx);
        \end{tikzpicture}
      \end{fitdiagram}
    \end{column}
  \end{columns}
\end{frame}

\begin{frame}{Tượng bị cản mắt}
  \begin{columns}[T,totalwidth=\textwidth]
    \begin{column}{0.42\textwidth}
      \begin{itemize}
        \item Tượng đi chéo hai ô.
        \item Ô giữa đường là mắt Tượng.
        \item Bị chặn mắt thì hướng đó không hợp lệ.
      \end{itemize}
      \begin{block}{Điểm khác cờ vua}
        Tượng cờ tướng không qua sông và bị blocker bởi mắt Tượng.
      \end{block}
    \end{column}
    \begin{column}{0.54\textwidth}
      \centering
      \includegraphics[height=0.68\textheight]{\figpath{figure-02b-elephant-eye.png}}
      \figcaption{3.b}{Tượng bị cản mắt.}
    \end{column}
  \end{columns}
\end{frame}

\begin{frame}{Tốt qua sông}
  \begin{columns}[T,totalwidth=\textwidth]
    \begin{column}{0.42\textwidth}
      \begin{itemize}
        \item Trước khi qua sông: Tốt chỉ tiến.
        \item Sau khi qua sông: có thêm nước ngang.
        \item Direction phụ thuộc bên Đỏ/Đen.
      \end{itemize}
    \end{column}
    \begin{column}{0.54\textwidth}
      \centering
      \includegraphics[height=0.68\textheight]{\figpath{figure-02c-pawn-river.png}}
      \figcaption{3.c}{Tốt đổi hướng sau khi qua sông.}
    \end{column}
  \end{columns}
\end{frame}

\begin{frame}{Pháo qua ngòi}
  \begin{columns}[T,totalwidth=\textwidth]
    \begin{column}{0.42\textwidth}
      \begin{itemize}
        \item Pháo đi như Xe trên ô trống.
        \item Muốn ăn quân phải có đúng một quân làm ngòi.
        \item Lingine tách quiet slides và leap captures.
      \end{itemize}
    \end{column}
    \begin{column}{0.54\textwidth}
      \centering
      \includegraphics[height=0.68\textheight]{\figpath{figure-03a-cannon-screen.png}}
      \figcaption{4.a}{Pháo cần đúng một ngòi để ăn.}
    \end{column}
  \end{columns}
\end{frame}

\begin{frame}{Xe quét rank/file}
  \begin{columns}[T,totalwidth=\textwidth]
    \begin{column}{0.42\textwidth}
      \begin{itemize}
        \item Xe di chuyển theo các tia ngang và dọc.
        \item Tập ô đi được phụ thuộc vào trạng thái chiếm đóng (occupancy) trên rank/file.
        \item Sử dụng lookup table để xác định nước đi hợp lệ với độ phức tạp $O(1)$.
      \end{itemize}
    \end{column}
    \begin{column}{0.54\textwidth}
      \centering
      \includegraphics[height=0.68\textheight]{\figpath{figure-03b-rook-rays.png}}
      \figcaption{4.b}{Xe quét rank/file.}
    \end{column}
  \end{columns}
\end{frame}

\begin{frame}{Lộ mặt Tướng}
  \begin{columns}[T,totalwidth=\textwidth]
    \begin{column}{0.42\textwidth}
      \begin{itemize}
        \item Hai Tướng không được đối mặt trên cùng file nếu không có quân chắn.
        \item Luật này ảnh hưởng trực tiếp tới \code{legal(m)}.
        \item Move hợp lệ phải đảm bảo Tướng không bị chiếu sau khi đi.
      \end{itemize}
    \end{column}
    \begin{column}{0.54\textwidth}
      \centering
      \includegraphics[height=0.68\textheight]{\figpath{figure-03c-flying-general.png}}
      \figcaption{4.c}{Lộ mặt Tướng.}
    \end{column}
  \end{columns}
\end{frame}

\begin{frame}{Backward attack scanner}
  \begin{columns}[T,totalwidth=\textwidth]
    \begin{column}{0.42\textwidth}
      Thay vì sinh toàn bộ nước đi của đối thủ:
      \begin{itemize}
        \item Lấy ô Tướng làm gốc.
        \item Quét các hướng có thể có attacker.
        \item Kiểm tra Tốt, Mã, Xe, Pháo, Tướng.
      \end{itemize}
    \end{column}
    \begin{column}{0.54\textwidth}
      \centering
      \includegraphics[height=0.68\textheight]{\figpath{figure-04a-backward-scanner.png}}
      \figcaption{4.d}{Quét ngược từ ô Tướng.}
    \end{column}
  \end{columns}
\end{frame}

\begin{frame}{Lọc move hợp lệ}
  \centering
  \begin{fitdiagram}
    \begin{tikzpicture}[node distance=0.55cm]
      \node[flowbox] (m) {pseudo move};
      \node[flowbox, below=of m] (do) {\code{do\_move}};
      \node[flowbox, below=of do] (check) {Tướng bị chiếu?};
      \node[goodbox, below left=0.65cm and 0.55cm of check] (legal) {giữ};
      \node[redbox, below right=0.65cm and 0.55cm of check] (drop) {loại};
      \draw[arrow] (m) -- (do);
      \draw[arrow] (do) -- (check);
      \draw[arrow] (check) -- (legal);
      \draw[arrow] (check) -- (drop);
    \end{tikzpicture}
  \end{fitdiagram}
\end{frame}

\begin{frame}{MoveList trên stack}
  \begin{columns}[T,totalwidth=\textwidth]
    \begin{column}{0.48\textwidth}
      \begin{block}{Thiết kế}
        \code{ArrayVec<Move, 128>}
      \end{block}
      \begin{itemize}
        \item Tránh cấp phát heap ở mỗi node.
        \item Dung lượng đủ cho branching factor cờ tướng.
        \item Phù hợp move ordering và search sâu.
      \end{itemize}
    \end{column}
    \begin{column}{0.48\textwidth}
      \centering
      \begin{fitdiagram}
        \begin{tikzpicture}[node distance=0.18cm]
          \foreach \i in {0,...,7} {
            \node[draw, minimum width=4.4cm, minimum height=0.42cm, fill=black!3] at (0,-0.45*\i) {\small Move slot \i};
          }
          \node[text=hustred, font=\bfseries] at (0,0.7) {stack buffer};
        \end{tikzpicture}
      \end{fitdiagram}
    \end{column}
  \end{columns}
\end{frame}



